%=====================================================================
\chapter{The Boundary with the Neighbouring Courses}\label{app:boundary}
%=====================================================================
\normalfigures

This course overlaps four others. The overlap is deliberate and bounded, and
this appendix states where the lines fall --- for students choosing options, for
colleagues designing a scheme of studies, and for anyone asking the fair
question of why this course exists at all.

\section{The Question This Appendix Answers}

\chref{what}'s Course Specification set out the position bluntly. HEC published
contents for \emph{Virtual Systems and Services} once, in 2017. The 2023
revision keeps the title without a specification. The 2025 revision
\textbf{removes the course}, dividing its ground between \emph{Cloud Computing},
a new major core, and \emph{Microservices Architecture and Docker Containers}.

The curriculum map that traces those editions reaches a blunt conclusion:
\emph{``Virtualisation is one chapter of cloud computing, not a peer of it.
Where both exist in a scheme of studies, retire Virtual Systems \& Services.''}

That is a reasonable position and this appendix does not dispute it in general.
What it does is state the specific case in which this course still earns three
credit hours, and what it must then teach that the others do not.

\begin{conceptbox}[The one-sentence answer]
A cloud course teaches you to \emph{use} the layer. This course teaches you what
is underneath it: which instruction trapped, which page table was walked twice,
which driver was split in half, and why the instance beside yours can make yours
slow.

\smallskip
That knowledge is what distinguishes someone who can run an estate from someone
who can only rent one. Where a department offers only \emph{Cloud Computing},
the material in Parts~I and~II of this handout is not taught anywhere in the
degree --- and it is the material that does not become obsolete.
\end{conceptbox}

\section{Cloud Computing}

The largest overlap, and the course HEC 2025 created to replace this one.

\rowcolors{2}{white}{lightbg}
\begin{longtable}{@{}L{4.6cm}L{5.0cm}L{5.0cm}@{}}
\toprule
\rowcolor{primary!15}
\textbf{Topic} & \textbf{Here} & \textbf{There} \\
\midrule
\endhead
Service models & One chapter (\chref{cloud}), as the commercial arrangement
around the preceding thirteen & Several weeks: the models, the providers, the
managed services, architecture patterns \\
Virtualization mechanism & \textbf{Parts I and II entire} --- Popek--Goldberg,
VT-x, EPT, the IOMMU, virtio, scheduling, overcommit & Usually one lecture,
descriptively \\
Cost & \chref{cloud}'s TCO method, and \chref{casestudies}'s migration that
cost more & Billing models, cost management tooling, FinOps practice \\
Storage and network & \chref{storage} and \chref{network}, as mechanisms you
could build & As cloud services you consume \\
Cloud-native architecture & Not taught & Microservices, serverless design,
event-driven patterns, the twelve factors \\
Managed services & Not taught & Databases, queues, identity, the provider's
catalogue \\
\bottomrule
\end{longtable}

\begin{alertbox}[If your department offers both, read this]
Running both courses in the same scheme of studies is defensible only if this
one keeps its centre of gravity in Parts~I--III and concedes Part~IV's cloud
material to the other.

\smallskip
Concretely: \chref{cloud} here should be taught as \emph{``the thing you have
been building, rented''} in a single session, not expanded. The material that
justifies this course's existence is \chref{vmm}, \chref{hardware},
\chref{cpumem} and \chref{paravirt} --- none of which a cloud course has room
for.

\smallskip
If a department can offer only one, and its graduates will mostly consume cloud
services rather than operate infrastructure, the curriculum map's advice is
right: teach \emph{Cloud Computing}.
\end{alertbox}

\section{Microservices Architecture and Docker Containers}

The other half of HEC 2025's replacement, and the overlap with Part~III.

\rowcolors{2}{white}{lightbg}
\begin{longtable}{@{}L{4.6cm}L{5.0cm}L{5.0cm}@{}}
\toprule
\rowcolor{primary!15}
\textbf{Topic} & \textbf{Here} & \textbf{There} \\
\midrule
\endhead
Namespaces, cgroups, overlayfs & \chref{containers} --- built by hand, as
kernel mechanisms & Usually assumed \\
Docker and images & \chref{docker} --- layers, digests, the build cache, the
OCI stack & The same ground, plus registries, CI integration, Compose \\
Kubernetes & \chref{kubernetes} --- the control loop, scheduling, requests and
limits, probes & Several weeks: operators, Helm, service mesh, GitOps,
autoscaling, stateful workloads \\
Microservice design & \textbf{Not taught} & The whole point of that course:
decomposition, API design, sagas, observability, team boundaries \\
Isolation runtimes & \chref{lightweight} --- microVMs, gVisor, Kata & Rarely \\
\bottomrule
\end{longtable}

The division is clean enough to state in one line: \textbf{this course teaches
what a container \emph{is}; that one teaches what to \emph{build} with
them.}

\section{System and Network Administration}

\begin{tight}
  \item \textbf{There:} users and permissions, backup \emph{policy}, directory
        services, patching processes, monitoring tooling, service management,
        the day-to-day operation of servers.
  \item \textbf{Here:} the same topics only where virtualization changes them
        --- \chref{management}'s fleet problems, \chref{migration}'s live
        migration and quorum, \chref{performance}'s host-side counters, which
        exist only because there is a hypervisor.
  \item \textbf{The test:} if the topic would be identical on physical servers,
        it belongs there.
\end{tight}

\section{CCNA and the Networking Track}

\begin{tight}
  \item \textbf{There:} switching, routing, subnetting, VLANs as a network
        technology, spanning tree, routing protocols, device configuration.
  \item \textbf{Here:} \chref{network} only --- the virtual switch, why VLANs
        run out at 4094, overlays and their MTU arithmetic, the control/data
        plane split, and where east--west traffic hides from physical
        monitoring.
  \item \textbf{The dependency:} this course \emph{assumes} the networking
        track's material. \chref{network} does not teach what a subnet is.
\end{tight}

\section{Data Science, Chapter 24}

The \emph{Data Science Fall 2026} handout has a chapter on cloud and edge
computing. The overlap is small and worth naming because students take both.

\begin{tight}
  \item \textbf{There:} where a model should run --- cloud or edge --- decided
        by bandwidth, latency, privacy and cost, with the arithmetic of
        telemetry volumes.
  \item \textbf{Here:} \chref{cloud} and \chref{casestudies}'s case 3, decided
        by the same kind of arithmetic applied to infrastructure rather than
        models.
  \item \textbf{Neither teaches the other.} A student who has taken both should
        notice that the method is identical and the quantities differ, which is
        a useful thing to notice.
\end{tight}

\section{What This Course Deliberately Does Not Teach}

Stated so that nobody is surprised, and so that a colleague designing a scheme
of studies knows what is left uncovered.

\rowcolors{2}{white}{lightbg}
\begin{longtable}{@{}L{5.0cm}L{9.8cm}@{}}
\toprule
\rowcolor{primary!15}
\textbf{Not taught here} & \textbf{Where it belongs} \\
\midrule
\endhead
Microservice decomposition and API design & \emph{Microservices Architecture
and Docker Containers} \\
Cloud-native application architecture, managed services, FinOps &
\emph{Cloud Computing} \\
Service mesh, Helm, GitOps, operators & \emph{Microservices}, or a DevOps
elective \\
Routing protocols, switching, subnetting & The networking track \\
Directory services, backup policy, user administration & \emph{System and
Network Administration} \\
Writing a hypervisor & A graduate operating-systems course. This book explains
how one works; it does not ask you to build one \\
Virtual \emph{reality} & Nowhere in this degree --- and it is what the 2017
outline's two set texts are about (Appendix~E) \\
\bottomrule
\end{longtable}

\begin{conceptbox}[A recommendation, stated plainly]
For a department deciding what to run:

\smallskip
\textbf{If you offer \emph{Cloud Computing} and \emph{Microservices}}, and have
room for only one more, this course is the one that supplies what neither does:
the mechanism beneath both. Teach it with its weight in Parts~I--III, and treat
Part~IV as revision.

\smallskip
\textbf{If you offer neither}, this course covers enough of both to be a
reasonable single elective, and Part~IV becomes load-bearing rather than
revision.

\smallskip
\textbf{If you offer \emph{Cloud Computing} only, and cannot add a second
course}, the curriculum map's advice stands --- but be aware that Parts~I
and~II then go untaught, and that is the part of this subject that was true in
1974 and will still be true in 2040.
\end{conceptbox}
